\chapter{Central Pattern Generators}
\label{ch:cpg}

\epigraph{%
  Rhythm is the most ancient form of motor control.
}{}

\section{CPGs in Biological Motor Control}

Central pattern generators are neural circuits that produce rhythmic motor output
without rhythmic sensory input or supraspinal commands. They underlie locomotion,
respiration, chewing, and other rhythmic behaviors.

\section{The Matsuoka Oscillator}

The most widely used CPG model in robotics is the Matsuoka oscillator
\cite{Matsuoka1985}:
\begin{align}
  \tau_1 \dot{x}_i &= -x_i - \beta v_i - \sum_{j \neq i} w_{ij} y_j + u_i + s_i, \\
  \tau_2 \dot{v}_i &= -v_i + y_i, \\
  y_i &= \max(0, x_i),
  \label{eq:ch8:matsuoka}
\end{align}
where $x_i$ is the membrane potential, $v_i$ is the adaptation variable, $y_i$ is the
output, $w_{ij}$ are inhibitory coupling weights, $u_i$ is the tonic input, and
$s_i$ is the sensory feedback.

\begin{lstlisting}[language=Python, caption={Matsuoka oscillator CPG implementation.}]
import numpy as np
from scipy.integrate import solve_ivp

class MatsuokaOscillator:
    """Two-neuron Matsuoka oscillator for CPG-based control."""

    def __init__(
        self,
        tau1: float = 0.1,
        tau2: float = 0.3,
        beta: float = 2.5,
        w_mutual: float = 2.0,
        tonic_input: float = 1.0
    ):
        self.tau1 = tau1
        self.tau2 = tau2
        self.beta = beta
        self.w = w_mutual
        self.u = tonic_input

    def dynamics(self, t: float, state: np.ndarray) -> np.ndarray:
        x1, v1, x2, v2 = state
        y1 = max(0.0, x1)
        y2 = max(0.0, x2)

        dx1 = (-x1 - self.beta * v1 - self.w * y2 + self.u) / self.tau1
        dv1 = (-v1 + y1) / self.tau2
        dx2 = (-x2 - self.beta * v2 - self.w * y1 + self.u) / self.tau1
        dv2 = (-v2 + y2) / self.tau2

        return np.array([dx1, dv1, dx2, dv2])

    def simulate(self, t_end: float = 5.0,
                 dt: float = 0.001) -> tuple[np.ndarray, np.ndarray]:
        state0 = np.array([0.1, 0.0, -0.1, 0.0])
        t_eval = np.arange(0, t_end, dt)
        sol = solve_ivp(self.dynamics, [0, t_end], state0,
                        t_eval=t_eval, method='RK45')

        # Extract outputs
        y1 = np.maximum(0, sol.y[0])
        y2 = np.maximum(0, sol.y[2])
        output = y1 - y2  # Net output

        return sol.t, output


# Simulate and print summary
cpg = MatsuokaOscillator(tau1=0.1, tau2=0.3)
t, output = cpg.simulate(t_end=5.0)

# Estimate frequency from zero crossings
crossings = np.where(np.diff(np.sign(output)))[0]
if len(crossings) >= 4:
    periods = np.diff(t[crossings[::2]])
    freq = 1.0 / np.mean(periods)
    print(f"CPG frequency: {freq:.2f} Hz")
    print(f"CPG amplitude: {np.max(np.abs(output[len(output)//2:])):.3f}")
\end{lstlisting}

\section{Phase Coupling and Inter-Limb Coordination}

Multiple CPGs can be coupled to produce coordinated multi-limb movement:
\begin{equation}
  \dot{\phi}_i = \omega_i + \sum_{j} K_{ij} \sin(\phi_j - \phi_i - \psi_{ij}),
  \label{eq:ch8:phase-coupling}
\end{equation}
where $\phi_i$ is the phase of the $i$-th oscillator, $\omega_i$ is its natural
frequency, $K_{ij}$ is the coupling strength, and $\psi_{ij}$ is the desired phase
relationship.

For quadruped locomotion:
\begin{itemize}
  \item Walk: $\psi = (0, \pi, \pi/2, 3\pi/2)$ (diagonal pairs alternating)
  \item Trot: $\psi = (0, \pi, \pi, 0)$ (diagonal pairs synchronized)
  \item Gallop: $\psi = (0, 0, \pi, \pi)$ (front-back alternating)
\end{itemize}

\section*{Exercises}

\begin{enumerate}
  \item \textbf{CPG Tuning.} Vary $\tau_1/\tau_2$ ratio and plot the resulting
        frequency and duty cycle. How does the ratio control the waveform shape?

  \item \textbf{Phase Coupling.} Implement 4 coupled oscillators and produce
        walk, trot, and gallop patterns by varying $\psi_{ij}$.

  \item \textbf{(Programming.)} Use the Matsuoka CPG to drive a simulated bipedal
        walker. Tune the CPG parameters for stable walking.
\end{enumerate}
